% Source: https://artofproblemsolving.com/community/c3h2391963_mathcounts_drills
% Stable identity: aops-topic-2391963-q14
Let the legs be $a$, $b$ and the hypotenuse $c$, so $a^2+b^2=c^2$ and $a+b+c=36$.

\begin{asy}
usepackage("amsmath");
size(4cm);
defaultpen(fontsize(10pt));
pair A=(0,0),B=(12,0),C=(0,9);
draw(A--B--C--cycle,linewidth(0.9));
draw((0,0.8)--(0.8,0.8)--(0.8,0));
label("$a$",(0,4.5),W); label("$b$",(6,0),S); label("$c$",(6,4.5),NE);
label("$a+b+c=36$",(4,-2));
\end{asy}

Substituting $c=36-a-b$:
\[a^2+b^2=(36-a-b)^2=1296-72a-72b+a^2+2ab+b^2,\]
so $2ab-72a-72b+1296=0$, which factors as
\[(36-a)(36-b)=648.\]
Both factors are positive integers (as $a,b<36$) and less than $36$. Writing $648=2^3\cdot3^4$, the factorisations of $648$ into two factors less than $36$ are only $24\cdot27$ (the next candidates $18\cdot36$ and $36\cdot18$ have a factor $36$). So $\{36-a,36-b\}=\{24,27\}$, i.e. $\{a,b\}=\{12,9\}$, and $c=15$.

The triangle is the $9$-$12$-$15$ right triangle, with area $\dfrac12\cdot9\cdot12=\boxed{54}$ square units.
